\emph{Översikt}
Den här modulen handlar om att låta datorn göra samma sak om och om igen.
Vi utgår från en uppgift som går att lösa genom att skriva ut samma steg
femtio gånger --- och ser varför man inte vill det.  Ur den kontrasten
växer tre sätt att uttrycka en upprepning en enda gång: \mintinline{python}{for},
som gör något för varje element i en samling; \mintinline{python}{while},
som upprepar tills ett villkor inte längre gäller; och \emph{rekursion},
där en funktion anropar sig själv på ett mindre problem.  Vi tittar
noggrant på vad som egentligen upprepas i en slinga, hur indraget avgör
vilka rader som hör till den, och hur en sträng kan gås igenom tecken för
tecken precis som vilken annan samling som helst.  Modulen avslutas med
att välja mellan de tre formerna på ett och samma problem, och med en
återanvändbar funktion för inmatning som frågar om igen tills användaren
har svarat begripligt.

\emph{Lärandemål}
Efter denna modul ska studenten kunna:
\begin{restatable}{lo}{IterLOFor}\label{IterLOFor}%
  Använda \mintinline{python}{for} för att göra samma sak en gång för
  varje element i en samling, och avgöra vilka rader som upprepas och
  vilka som inte gör det. % Lo04
\end{restatable}
\begin{restatable}{lo}{IterLOWhile}\label{IterLOWhile}%
  Använda \mintinline{python}{while} för att upprepa så länge ett villkor
  gäller, och se till att villkoret kan bli falskt. % Lo04
\end{restatable}
\begin{restatable}{lo}{IterLORecursion}\label{IterLORecursion}%
  Läsa och skriva rekursiva funktioner med basfall och rekursivt anrop på
  ett mindre problem. % Lo04
\end{restatable}
\begin{restatable}{lo}{IterLOChoose}\label{IterLOChoose}%
  Välja mellan \mintinline{python}{for}, \mintinline{python}{while} och
  rekursion utifrån hur problemet ser ut, och motivera valet. % Lo04, Lo02
\end{restatable}
\begin{restatable}{lo}{IterLOString}\label{IterLOString}%
  Se en sträng som en följd av tecken som går att gå igenom på samma sätt
  som vilken annan samling som helst. % Lo04
\end{restatable}
\begin{restatable}{lo}{IterLOInput}\label{IterLOInput}%
  Skriva en återanvändbar funktion för inmatning som frågar om igen tills
  användaren har matat in data som går att använda. % Lo08, Lo01, Lo03
\end{restatable}

\emph{Förkunskaper}
Modulen förutsätter veckorna om funktioner och variabler samt om inmatning
och felhantering: studenten ska kunna skriva en funktion med parametrar och
returvärde, läsa in text med \mintinline{python}{input}, omvandla den med
\mintinline{python}{int} och fånga ett särfall med
\mintinline{python}{try}/\mintinline{python}{except}, samt använda
\mintinline{python}{if}, \mintinline{python}{elif} och
\mintinline{python}{else}.

\emph{Läsning}
Kompletterande material finns i kursens videogenomgång \emph{Upprepningar}.
Listor och tupler --- de samlingar en \mintinline{python}{for}-slinga oftast
går igenom --- behandlas i modulens egna material \emph{Listor} och
\emph{Tupler}.
